% Copyright 2026 Imanol Felix Supo Mamani. Licensed under LPPL 1.3c.
% See NOTICE.md and MANIFEST.txt for work scope and maintenance.

\documentclass[11pt]{article}
\usepackage[margin=25mm]{geometry}
\usepackage[T1]{fontenc}
\usepackage{ormath}
\usepackage{array,longtable}
\usepackage[hidelinks]{hyperref}
\hypersetup{pdftitle={ORmath user guide: experimental version 0.1.0},
  pdfauthor={Imanol Felix Supo Mamani},
  pdfsubject={Typesetting mathematical optimization models and notation}}
\setlength{\parindent}{0pt}
\setlength{\parskip}{4pt}
\raggedbottom
\newcommand{\key}[1]{\texttt{#1}}
\newenvironment{sourceexample}{\par\noindent\begin{minipage}{\linewidth}\small}{\end{minipage}\par}
\newcommand{\responsiveexample}{\InputIfFileExists{examples/fixtures/responsive-model.tex}{}{\input{responsive-model.tex}}}
\title{ORmath\\[4pt]
  {\large Typesetting mathematical optimization models and notation}}
\author{Imanol Felix Supo Mamani}
\date{Experimental version 0.1.0, September 20, 2026}
\makeatletter
\renewcommand{\maketitle}{\begin{center}
  {\Large\@title\par}\vspace{2pt}
  {\normalsize\@author\par}\vspace{1pt}
  {\small\@date\par}
\end{center}}
\makeatother
\begin{document}
\maketitle
\section{What is ORmath?}
ORmath is a lightweight, width-aware toolkit for mathematical optimization
models. Use ordinary LaTeX mathematics; let the package arrange notation,
index domains, and equation numbers. It typesets models without solving or
validating them. The API is provisional; start with the lightweight workflow below.
Load the package with \verb|\usepackage{ormath}|.

\section{Quick start}
Copy \key{ormath.sty} and \key{ormath-input.code.tex} beside your document.
Save this complete example as \key{00-quickstart.tex} and compile twice.
Here $I$ contains $n$ products, $c_i$ is unit profit, and $u_i\geq0$ is capacity.
\begin{sourceexample}
\begin{verbatim}
\documentclass{article}
\usepackage{ormath}
\begin{document}
\begin{ornotation}
  \set[size=n]{I}{Products}
  \param{c_i}{Unit profit}\for{i in I}
  \param{u_i}{Capacity}\for{i in I}
  \var{x_i >= 0}{Production}\for{i in I}
\end{ornotation}
\begin{ormodel}
  \maximize \sum_{i\in I} c_i x_i \\
  \st x_i <= u_i \for{i in I} \label{capacity}
\end{ormodel}
Capacity is respected by constraint~\eqref{capacity}.
\end{document}
\end{verbatim}
\end{sourceexample}
\textbf{Output.} Each product's production is bounded by its capacity.
\begin{ornotation}
  \set[size=n]{I}{Products}
  \param{c_i}{Unit profit}\for{i in I}
  \param{u_i}{Capacity}\for{i in I}
  \var{x_i >= 0}{Production}\for{i in I}
\end{ornotation}
\begin{ormodel}
  \maximize \sum_{i\in I} c_i x_i \\
  \st x_i <= u_i \for{i in I} \label{capacity}
\end{ormodel}
Capacity is respected by constraint~\eqref{capacity}.

\clearpage
\section{The recommended workflow}
Use \key{ornotation} for definitions, then \key{ormodel} for the objective and
constraints. Short commands are local to those environments. No layout options
are needed for the default workflow.

\subsection{Extend the notation}
Continuing the quick start, let $K$ contain $m$ machines. A defined set can hold
pairs rather than consecutive indices; give its mathematical definition directly.
\begin{sourceexample}
\begin{verbatim}
\begin{ornotation}
  \set[size=m]{K}{Machines}
  \setdef{A}{(i,k)\in I\times K}{Product-machine pairs}
  \param{a_{ik}}{Processing time}\for{(i,k) in A}
  \var{z_{ik}\in\{0,1\}}{Assignment}\for{(i,k) in A}
\end{ornotation}
\end{verbatim}
\end{sourceexample}
\begin{ornotation}
  \set[size=m]{K}{Machines}
  \setdef{A}{(i,k)\in I\times K}{Product-machine pairs}
  \param{a_{ik}}{Processing time}\for{(i,k) in A}
  \var{z_{ik}\in\{0,1\}}{Assignment}\for{(i,k) in A}
\end{ornotation}

\subsection{Choose the objective sense and write rows}
Minimizing negative profit gives the same optimum as maximizing profit. Use
\verb|\minimize| with a minus sign; the capacity constraint stays unchanged.
\begin{sourceexample}
\begin{verbatim}
\begin{ormodel}
  \minimize -\sum_{i\in I}c_i x_i \\
  \st x_i <= u_i \for{i in I} \label{capacity-min}
\end{ormodel}
Capacity constraint:~\eqref{capacity-min}.
\end{verbatim}
\end{sourceexample}
\begin{ormodel}
  \minimize -\sum_{i\in I}c_i x_i \\
  \st x_i <= u_i \for{i in I} \label{capacity-min}
\end{ormodel}
Capacity constraint:~\eqref{capacity-min}.

\subsection{The commands to remember}
\par\noindent\begin{tabular}{@{}l>{\raggedright\arraybackslash}p{.63\linewidth}@{}}
\verb|\set| / \verb|\setdef| & Name a set or give its definition.\\
\verb|\param| / \verb|\var| & Describe a parameter or complete variable declaration.\\
\verb|\for| & Attach an index domain, such as \verb|i in I|.\\
\verb|\minimize| / \verb|\maximize| & Begin the objective.\\
\verb|\st| & Introduce the constraints; \verb|\\| separates rows.
\end{tabular}\par
Use ordinary LaTeX for mathematics. This is the complete lightweight workflow;
customization and the structured API are optional.

\clearpage
\section{Compact notation}
\subsection{Sets}
Inside \key{ornotation}, use \verb|\set[options]{symbol}{description}|.
A plain set preserves your symbol without inventing members. Select at most one
representation: \key{size}, \key{elements}, or \key{definition} (all empty by default).
\verb|\setdef{symbol}{definition}{description}| is a direct defined-set form.
\begin{sourceexample}
\begin{verbatim}
\begin{ornotation}
  \set{C}{Customers with arbitrary identifiers}
  \set[size=n]{I}{Consecutive product indices}
  \set[elements={A,B,C}]{H}{Harvest types}
  \setdef{M}{(i,t)\in I\times T}{Product-period pairs}
\end{ornotation}
\end{verbatim}
\end{sourceexample}
Here $T$ denotes the planning-period set; define it in your application's notation.
\key{size=|C|} explicitly requests an enumeration ending at $|C|$.

\subsection{Parameters and variables}
Use \verb|\param{math}{description}| and \verb|\var{math}{description}|.
Descriptions are ordinary text and can contain inline mathematics. Group related
symbols when one description is clearer. A variable is the complete declaration
you supply: bounds and domains are explicit, with no implicit variable type.
\begin{sourceexample}
\begin{verbatim}
\begin{ornotation}
  \set[size=H]{T}{Periods}
  \param{c_{it},f_{it},h_{it}}{Production, setup, and holding costs}
    \for{i in I,t in T}
  \var{x_{it} >= 0}{Production}\for{i in I,t in T}
  \var{y_{it}\in\{0,1\}}{Setup indicator}\for{i in I,t in T}
\end{ornotation}
\end{verbatim}
\end{sourceexample}

\subsection{Index domains with \texorpdfstring{\texttt{\string\for}}{for}}
Attach \verb|\for{i in I,t in T}| to the preceding declaration or model row.
Native membership, such as \verb|\for{i\in I,t\in T}|, works too.
Tuple memberships are supported. For arbitrary mathematical conditions, use
\verb|\for*{t>0}|; the starred form renders raw mathematics without membership parsing.

\subsection{Controlled shorthand}
\par\noindent\begin{tabular}{@{}ll@{}}
\textbf{Input} & \textbf{Rendered meaning}\\
\verb|x >= 0| & $x\geq0$\\
\verb|x <= U| & $x\leq U$\\
\verb|\for{i in I}| & $\forall\,i\in I$\\
\verb|\for{i in I,j in J}| & $\forall\,i\in I,j\in J$\\
\verb|\for{(i,j) in A}| & $\forall\,(i,j)\in A$
\end{tabular}\par
Inequality shorthand applies at the top level in raw variables and lightweight
model rows; \key{in} applies only inside unstarred \verb|\for|. It is local and
preserves ordinary LaTeX. Write native operators inside braced groups or macro
arguments. There is no \verb|!=| or \verb|==| shorthand: use \verb|\neq| and
\verb|=|. In particular, \verb|n!=k| keeps factorial followed by equality.

\clearpage
\section{Models}
\subsection{Objectives and constraints}
For a minimum-cost variant, let $p_i\geq0$ be unit cost and $d_i\geq0$
required production. Keep $x_i\geq0$ as in the quick start.
Begin literally with \verb|\minimize| or \verb|\maximize|. Exactly one objective
is required. Put it first, then the constraints. Use explicit \verb|\\| between
rows; blank lines are whitespace. \verb|\st| introduces constraints;
\verb|\subjectto| is its longer synonym. \verb|\for| supplies the index domain.
\begin{sourceexample}
\begin{verbatim}
\begin{ormodel}
  \minimize \sum_{i\in I}p_i x_i \label{cost} \\
  \st x_i >= d_i \for{i in I} \label{minimum} \\
  x_i >= 0 \for{i in I} \notag
\end{ormodel}
\end{verbatim}
\end{sourceexample}
\begin{ormodel}
  \minimize \sum_{i\in I}p_i x_i \label{cost} \\
  \st x_i >= d_i \for{i in I} \label{minimum} \\
  x_i >= 0 \for{i in I} \notag
\end{ormodel}

\subsection{Labels and references}
Use ordinary top-level \verb|\label|, then \verb|\ref| or \verb|\eqref| in text.
The objective is~\eqref{cost}; the minimum-production constraint is~\eqref{minimum}.
Compile twice. \verb|\notag| suppresses one row's number and counter increment;
that row cannot carry a label. Definitions in \key{ornotation} have no equation
number and cannot carry labels.

\par\noindent\begin{tabular}{@{}l>{\raggedright\arraybackslash}p{.61\linewidth}@{}}
\key{numbering=rows} & Default: one number per numbered row; row labels allowed.\\
\key{numbering=model} & One number beside the first row; use only an environment label.\\
\key{numbering=none} & No counter increment; labels are rejected.
\end{tabular}\par
With row numbering, an environment label refers to the first numbered row.
Do not combine \verb|\notag| with \key{numbering=model}.

\subsection{Long expressions and manual breaks}
Protect an internally multiline expression with one outer brace group. The
internal \verb|&| and \verb|\\| then reach amsmath unchanged; the following
outer \verb|\\| still separates model rows. These are authored breaks.
\begin{sourceexample}
\begin{verbatim}
\begin{ormodel}
  \minimize {&\sum_{i\in I}p_i x_i\\
             &{}+\sum_{k\in K}f_k y_k} \\
  \st x_i >= 0 \for{i in I}
\end{ormodel}
\end{verbatim}
\end{sourceexample}
Here $f_k,y_k$ are machine setup costs and indicators. Use \verb|{}+| or
\verb|{}-| to preserve binary-operator spacing on a continuation line.
ORmath does not automatically split arbitrary mathematical expressions.

\clearpage
\subsection{Indexed summations}
Inside \texttt{ormodel}, the optional \verb|\Sum| shorthand produces an indexed
operator and leaves all following mathematics untouched. One command always
means one summation symbol. Comma-separated memberships share its subscript;
write two commands explicitly for nested summation symbols.
\begin{sourceexample}
\begin{verbatim}
\begin{ormodel}[numbering=none]
  \minimize \Sum{i in I,j in J_i} c_{ij}x_{ij} \\
  \st \Sum{i in I}\Sum{j in J_i} c_{ij}x_{ij} <= B
\end{ormodel}
\end{verbatim}
\end{sourceexample}
\begin{ormodel}[numbering=none]
  \minimize \Sum{i in I,j in J_i} c_{ij}x_{ij} \\
  \st \Sum{i in I}\Sum{j in J_i} c_{ij}x_{ij} <= B
\end{ormodel}
Use \verb|\Sum{(i,j) in A}| for tuple membership and
\verb|\Sum{i in \mathcal{I}}| for a set written with ordinary math macros.
Brace groups and commas inside parentheses are protected. Set expressions and
macros are preserved without expansion by the parser.

An optional single \texttt{where} follows the complete membership list. Its
condition is ordinary mathematical LaTeX; condition commas are not membership
separators. Only top-level \verb|>=| and \verb|<=| are normalized, following
existing input policy. Neither \verb|!=| nor \verb|==| is rewritten, and words
such as \texttt{and} or \texttt{or} have no special meaning. Write
\verb|\land|, \verb|\lor|, \verb|\neq| and \verb|\in| normally in conditions.

Standard LaTeX remains available:
\begin{sourceexample}
\begin{verbatim}
\[
  \sum_{\substack{i\in I,\;j\in J\\a_{ij}>0}}
    c_{ij}x_{ij}
\]
\end{verbatim}
\end{sourceexample}
The optional convenience form uses a colon before the condition:
\begin{sourceexample}
\begin{verbatim}
\begin{ormodel}[numbering=none]
  \minimize \Sum{i in I,j in J where a_{ij}>0}
    c_{ij}x_{ij}
\end{ormodel}
\end{verbatim}
\end{sourceexample}
\begin{ormodel}[numbering=none]
  \minimize \Sum{i in I,j in J where a_{ij}>0} c_{ij}x_{ij}
\end{ormodel}

The public namespaced command \verb|\orsum| works in ordinary math and in
structured objectives and constraints.

The short alias \verb|\Sum| exists only inside \texttt{ormodel}; any previous
meaning is restored afterward. Neither command changes standard \verb|\sum|.

The starred form supplies its subscript verbatim,
bypassing membership parsing and operator normalization:
\begin{sourceexample}
\begin{verbatim}
\[
  \orsum*{\substack{i\in I\\j\in J_i\\(i,j)\notin A}}
    x_{ij}
\]
\end{verbatim}
\end{sourceexample}
Use \verb|\Sum*| inside a model for the same escape. Each convenience membership
needs an index and a set. Empty conditions, repeated \texttt{where}, and literal
\texttt{in} memberships after \texttt{where} produce actionable errors suggesting
the raw form. Keywords are literal and space-delimited; macros cannot supply
hidden grammar. Advanced subscripts, custom punctuation and multiline conditions
belong in the starred form or standard LaTeX.

The following realistic example uses the same source as the visual fixture:
\begin{sourceexample}
\begin{verbatim}
\begin{ormodel}
  \minimize \Sum{i in I} L_i
    + \Sum{i in I,j in J} c_{ij}x_{ij} \\
  \st \Sum{j in J} a_{ij}x_j <= b_i \for{i in I} \\
  \Sum{(i,j) in A where c_{ij}>0} y_{ij} <= M
\end{ormodel}
\end{verbatim}
\end{sourceexample}
\input{examples/fixtures/summation-model.tex}

\section{Responsive layout}
The same source adapts to the surrounding \verb|\linewidth|. This feasibility
model has activities $J$, resource constraints $I$, periods $T$, resource usage
$a_{ij}$, and capacities $b_{it}$. The constant objective asks only for feasibility.
\begin{sourceexample}
\begin{verbatim}
\begin{ormodel}[numbering=none]
  \minimize 0 \\
  \st \sum_{j\in J}a_{ij}x_{jt} <= b_{it}
    \for{i in I,t in T} \\
  x_{jt} >= 0 \for{j in J,t in T}
\end{ormodel}
\end{verbatim}
\end{sourceexample}
\textbf{Wide output}
\responsiveexample
\textbf{Narrow output (145pt local width)}
\par\noindent\begin{minipage}{145pt}
\responsiveexample
\end{minipage}\par
Both outputs read the identical fixture. The long capacity row needs its domain
below; the shorter nonnegativity row can retain its domain inline. Neither the
mathematics nor the font size changes.

\subsection{What the defaults do}
Notation and models begin at the local left edge. Each definition keeps its
mathematical declaration, colon, description, and domain on one line when they
fit. Safe gaps decrease before long descriptions wrap with hanging indentation;
domains move below only when necessary. Short rows stay compact beside long rows.
No automatic math-font styling, font shrinking, or margin invasion occurs.
User punctuation remains literal; model rows with domains default to a semicolon.

\subsection{Where width adaptation stops}
The model renderer measures expressions and index domains and reserves marker
and equation-number space. It adapts component placement, without inferring
mathematics or scaling boxes. Long unbreakable expressions still need authored
breaks. A multiline row is indivisible and cannot split across pages.

\subsection{Local model domains}
Every model row places its domain immediately after its own expression.
Equation tags remain separate from domains and use the adaptive placement
described below. With \key{domain-gap=auto},
the preferred and maximum gap is 0.8em and the minimum is 0.3em, measured in
the row's selected font. The renderer reserves marker and tag space, measures
the expression (including any separator) and domain, and uses the remaining
width to reduce the gap. If they still cannot fit, only that domain moves below.
A longer neighboring row never changes the gap or placement of a short row.

Use \key{domain-gap=1em} or \key{domain-gap=10pt} for a fixed model gap.
\key{domain-position=right} forces inline placement and can overflow;
\key{domain-position=below} forces a continuation line. No two-column switch,
font shrinking, or automatic margin extension is used. \key{layout=wide}
retains its existing forced-inline meaning; stacked layout takes precedence.

The default is \key{domain-separator=\{;\}}, following the author's editorial
preference.\par
Select \key{domain-separator=\{,\}} for commas.\par
Select \key{domain-separator=\{\}} for none. The renderer appends punctuation
only to expressions with domains, including when the domain moves below.
A literal trailing comma, semicolon, or period takes precedence, avoiding doubled
punctuation. Punctuation hidden in a macro or behind a spacing command is not
inferred; in that case author the punctuation and leave the separator empty.
The separator applies to models; notation retains its colon/prose conventions.

\clearpage
\subsection{Compact responsive equation tags}
The defaults are \key{tag-gap=auto} and \key{tag-position=auto}.
Available \verb|\linewidth| limits a model; it is not a width the model must fill.
Objectives and constraints share a compact tag anchor when their natural widths
permit it. A much longer row extends its own tag locally instead of pushing
every short row toward the page edge. No class names or page geometry are tested.

\begin{sourceexample}
\begin{verbatim}
\begin{ormodel}
  \minimize \sum_i c_i x_i \\
  \st x_i <= 10 \for{i in I} \\
  \sum_j a_{ij}x_j <= b_i \for{i in I} \\
  y_i+z_i <= 1 \for{i in I}
\end{ormodel}
\end{verbatim}
\end{sourceexample}
\input{examples/fixtures/tag-target.tex}

For each row, the renderer measures the marker, expression, domain and actual
upcoming tag in the selected font. Auto tag separation prefers one twenty-fifth
of the available width after markers, bounded between 0.75em and 2em. It can
compress to 0.5em when needed. Auto domain and tag gaps compress before a domain
moves below. The tag then remains on the expression's first line, and the
domain continuation can use the width below it. Mathematics is not shrunk.

Natural row extents determine a shared anchor with a bounded extension from
the central row width. Exceptional rows keep a local tag. Shared alignment
adds padding to shorter rows, but never consumes arbitrary spare page width.
The implementation and outlier rule are documented in ARCHITECTURE.md.

Use \key{tag-gap=0.8em} for a tighter exact gap.\par
Use \key{tag-gap=2em} for more breathing room, or \key{tag-gap=12pt} for a
dimension. An explicit gap places each tag that exact distance after its own
first-line content; tags may consequently have different horizontal positions.
These options apply at global setup, named-style, model and structured-row levels.

Use \key{tag-position=right} when a publisher requires conventional outer-right
tags. This mode takes precedence over the compact anchor; \key{tag-gap} controls
the minimum fitting reservation rather than the distance to the right edge.
Forced/manual choices can overflow an extremely narrow context.
The existing \key{numbering=rows}, \key{numbering=model}, \key{numbering=none},
ordinary labels/references and \verb|\notag| retain their counter semantics.

\clearpage
\section{Customization and styles}
Begin with the defaults. For repeated visual choices, define a style and select
it on an environment. This complete fragment uses the same lightweight rows.
\begin{sourceexample}
\begin{verbatim}
\ormathstyle{paper}{density=compact,row-gap=1.5pt}
\begin{ormodel}[style=paper]
  \minimize \sum_{i\in I}p_i x_i \\
  \st x_i >= d_i \for{i in I}
\end{ormodel}
\end{verbatim}
\end{sourceexample}
\verb|\ormathsetup{density=normal}| changes document defaults. Setup and style
definitions follow TeX grouping: top-level definitions persist, grouped ones
remain local. Styles contain visual keys and cannot select another \key{style}.
The precedence is element options, environment options, named style, global
setup, then defaults. Within a layer, explicit spacing overrides density,
regardless of key order.

\subsection{Option reference}
Defaults below describe the compact preset. All environments accept visual keys,
\key{style}, and \key{title}. Model-only keys are identified explicitly.
Dimensions must be nonnegative TeX dimensions.
\begingroup\small
\setlength{\LTleft}{0pt}\setlength{\LTright}{\fill}
\begin{longtable}{@{}>{\raggedright\arraybackslash}p{.21\linewidth}>{\raggedright\arraybackslash}p{.25\linewidth}>{\raggedright\arraybackslash}p{.13\linewidth}>{\raggedright\arraybackslash}p{.33\linewidth}@{}}
\textbf{Option} & \textbf{Accepted values} & \textbf{Default} & \textbf{Meaning}\\\hline
\endfirsthead
\textbf{Option} & \textbf{Accepted values} & \textbf{Default} & \textbf{Meaning}\\\hline
\endhead
\key{layout} & auto, wide, stacked & auto & Automatic placement or explicit component layout.\\
\key{domain-position} & auto, right, below & auto & Domain placement; forced right can overflow a narrow model.\\
\key{font-size} & auto, normal, small, footnotesize, scriptsize & auto & Inherit the font size or explicitly select one.\\
\key{density} & relaxed, normal, compact, dense & compact & Set row and outer spacing; automatic notation gap.\\
\key{row-gap} & dimension & 2pt & Additional space after each row.\\
\key{space-before} & dimension & 4pt & Vertical space before a block.\\
\key{space-after} & dimension & 4pt & Vertical space after a block.\\
\key{domain-gap} & auto or dimension & auto & Bounded model gap; density-dependent notation gap.\\
\key{domain-separator} & math tokens & ; & Model punctuation before a domain.\\
\key{tag-gap} & auto or dimension & auto & Model tag separation; explicit dimensions are row-local.\\
\key{tag-position} & auto, right & auto & Compact responsive model anchor or conventional outer-right tags.\\
\key{extend-left} & dimension & 0pt & Explicit width into the left margin.\\
\key{extend-right} & dimension & 0pt & Explicit width into the right margin.\\
\key{style} & defined style name & empty & Apply a named visual style.\\
\key{title} & text & empty & Add a block title.\\
\key{numbering} & rows, model, none & rows & Model equation-number policy.\\
\key{label} & label key & empty & Model environment reference.\\
\key{subject-to} & text & s.t. & Model constraint marker; e.g. \key{\{s.a.\}}.
\end{longtable}
\endgroup

\clearpage
\subsection{Spacing presets and placement details}
These values are reference information; ordinary writing needs no spacing setup.
\par\noindent\begin{tabular}{@{}lrrrr@{}}
\textbf{Density} & \textbf{Row gap} & \textbf{Before} & \textbf{After} & \textbf{Domain gap}\\
relaxed & 7pt & 12pt & 12pt & 16pt\\
normal & 4pt & 8pt & 8pt & 12pt\\
compact & 2pt & 4pt & 4pt & 8pt\\
dense & 1pt & 2pt & 2pt & 6pt
\end{tabular}\par
The last column describes automatic notation gaps. All presets now retain
\key{domain-gap=auto}; models use the bounded relative gap described above.
In notation, the default 8pt prose-domain gap may decrease to 3pt after a line
fails to fit. The 3pt colon-prose gap may decrease to 2pt. This occurs before
wrapping, without shrinking the mathematical font.

\begin{description}
\item[\key{layout=wide}] Model components stay side by side; notation prose
still wraps when needed.
\item[\key{layout=stacked}] Model domains go below their expressions;
notation prose goes below its mathematical prefix.
\item[\key{domain-position}] In auto/wide models, right forces a side domain;
stacked layout takes precedence. In notation, auto/right flow with prose when
width permits; below places the domain on a continuation line after the description.
\end{description}

\subsection{Element-level styles}
Element-level control is available through the advanced structured API in the
next section. Its row keys are \key{font-size}, \key{row-gap}, \key{domain-gap},
\key{domain-separator}, \key{domain-position}, \key{tag-gap}, and
\key{tag-position}. A row style may contain only those keys.
It overrides environment values, then explicit element options override it.
Styles never change mathematical notation or objective sense.

\clearpage
\section{Advanced structured API}
The lightweight API is recommended for ordinary model writing. The structured
API supplies additional metadata and element-level visual control, useful for
advanced customization or assembled rows. Both feed the same semantic records
and renderer where applicable; they are not separate layout engines.

\key{orsets}, \key{orparameters}, and \key{orvariables} hold structured
\verb|\set|, \verb|\parameter|, and \verb|\variable| declarations. Model rows use
\verb|\objective| and \verb|\constraint|. Keep one model input style per environment.
This example uses the minimum-cost variant from the models section.
\begin{sourceexample}
\begin{verbatim}
\begin{orsets}
  \set[size=n]{I}{Products}
\end{orsets}
\begin{orparameters}
  \parameter[for={i\in I}]{p_i}{Unit cost}
  \parameter[for={i\in I}]{d_i}{Required production}
\end{orparameters}
\begin{orvariables}
  \variable[type=nonnegative,for={i\in I}]{x_i}{Production}
\end{orvariables}
\begin{ormodel}
  \objective[min]{\sum_{i\in I}p_i x_i}
  \constraint[for={i\in I},font-size=small]{x_i\geq d_i}
\end{ormodel}
\end{verbatim}
\end{sourceexample}
\begin{orsets}
  \set[size=n]{I}{Products}
\end{orsets}
\begin{orparameters}
  \parameter[for={i\in I}]{p_i}{Unit cost}
  \parameter[for={i\in I}]{d_i}{Required production}
\end{orparameters}
\begin{orvariables}
  \variable[type=nonnegative,for={i\in I}]{x_i}{Production}
\end{orvariables}
\begin{ormodel}
  \objective[min]{\sum_{i\in I}p_i x_i}
  \constraint[for={i\in I},font-size=small]{x_i\geq d_i}
\end{ormodel}

\clearpage
\subsection{Structured metadata reference}
\par\noindent\begin{tabular}{@{}l>{\raggedright\arraybackslash}p{.7\linewidth}@{}}
\key{for} & Native mathematical index domain; empty by default.\\
\key{label} & Model-row reference; empty by default, subject to numbering policy.\\
\key{type} & Variable: continuous (default), nonnegative, integer, binary.\\
\key{min} / \key{max} & Objective sense; min is the default.\\
\key{style} & Named row style, plus explicit row keys from customization.
\end{tabular}\par
Structured variable types append a mathematical domain. Continuous and integer
variables are unrestricted in sign; add nonnegativity explicitly if needed.
The inherited upright $\mathrm R$ and $\mathrm Z$ letters require no extra math
font package. Raw lightweight \verb|\var| never appends a type.

Set representation keys \key{size}, \key{elements}, and \key{definition} have the
same meanings as in compact notation. Model objectives and constraints accept
\key{for}, \key{label}, and row visual keys. For example, metadata after the sense
can be written as \verb|\objective[max,label=obj:profit]{c^T x}| inside a structured
model. Labels on notation declarations are rejected.

\subsection{Namespaced forms and scope}
Namespaced structured commands are \verb|\orset|, \verb|\orparameter|,
\verb|\orvariable|, \verb|\orobjective|, and \verb|\orconstraint|. Lightweight
adapters also expose \verb|\orvar|, \verb|\orsetdef|, and \verb|\orfor| in their
checked contexts. The short aliases exist only within their corresponding
ORmath environments; prior meanings are restored afterwards. These declaration
commands diagnose use outside the required context. The mathematical operator
\verb|\orsum| also works in ordinary math, with \verb|\Sum| local to models.

\subsection{Package loading}
Load ORmath with \verb|\usepackage{ormath}|. Copy \key{ormath.sty} and
\key{ormath-input.code.tex} together for a local installation. No legacy or
compatibility loaders are provided. No claim is made about CTAN availability.

\clearpage
\section{Realistic example}
This optional production-routing formulation combines four cost terms, inventory
balances, conditional arc sums, shared capacity, and integer/binary decisions.
It uses the advanced structured API just described. The repository also has
lightweight planning examples and an API comparison using the same renderer.

The set $\mathcal I$ indexes sites and destinations, $\mathcal T$ periods,
$\mathcal M=\mathcal I\times\mathcal T$ site-period pairs, and $\mathcal A$ feasible
shipment arcs by period. Variables $x,s,z,y$ represent production, inventory,
shipment quantity, and setup. Parameters $c,f,q,h$ are costs; $d,U,B,a,D,S$ and
the initial/target stock values specify demands and capacities. Full definitions
are in \key{examples/fixtures/planning-data.tex}.
\par\noindent\begin{minipage}{.65\linewidth}
\InputIfFileExists{examples/fixtures/planning-model.tex}{}{\input{planning-model.tex}}
\end{minipage}\par
The shared body is used in \key{05-long-model.tex},\newline
\key{06-two-column.tex}, and \key{07-narrow-layout.tex}. These tested widths
do not certify every formulation or class.

\clearpage
\section{Engine compatibility}
The example and regression suite has been executed under pdfLaTeX, XeLaTeX,
and LuaLaTeX on Windows/MiKTeX. Unicode-engine smoke tests use fontspec and
unicode-math with Latin Modern fonts. Inherited fonts and accented prose are
covered; all fonts, publisher classes, and operating systems are not certified.
A LaTeX kernel dated 2022-06-01 or newer is required; amsmath is the only
external package dependency of the toolkit. The guide additionally uses
geometry, array, longtable, fontenc, and hyperref.

\section{Limitations and development status}
Version 0.1.0 is experimental. ORmath does not parse arbitrary mathematics,
infer set membership, solve models, or validate their mathematical meaning.
Explicit breaks are needed for some long expressions. Nested ORmath
environments, verbatim in mathematical content, and macro-generated row structure
are unsupported. Models align expression starts rather than arbitrary relations.
Very narrow contexts can overflow and are diagnosed.

An internally multiline row cannot split across pages. Continuation titles and
subject-to markers are not repeated on later pages. Measurement can evaluate
material more than once: avoid global side effects, counter increments, and
labels hidden inside mathematical groups. Use top-level lightweight labels or
structured label keys. Braced groups are opaque to shorthand normalization;
optional/unbraced macro arguments cannot be inferred, so use native operators
or a protective group for literal content.

\subsection{Copyright and license}
Copyright 2026 Imanol Felix Supo Mamani. Licensed under the LaTeX Project Public
License, version 1.3c. LPPL maintenance status: author-maintained.
Current Maintainer: Imanol Felix Supo Mamani. See LICENSE, NOTICE.md, and
MANIFEST.txt for the grant and work's scope. No public release has occurred.

\subsection{Reproduce the checks}
From the repository root:
\begin{sourceexample}
\begin{verbatim}
l3build check
python scripts/compile.py --engine all
l3build doc
python scripts/check_manual.py --engine all
\end{verbatim}
\end{sourceexample}
The last command extracts and compiles the executable guide examples. Build
commands above are shell commands, not LaTeX source. Actual results and
verification limits are recorded in \key{TEST-RESULTS.md}.
\end{document}
